\section{Metode Pendekatan}

Penelitian ini menggunakan pendekatan kuantitatif. Penelitian kuantitatif menguji teori secara objektif dengan memeriksa hubungan antarvariabel; variabel tersebut diukur dengan instrumen sehingga data berupa angka dapat dianalisis secara statistik \parencite{creswell2018}. Data penelitian ini berupa jumlah dan laju penandaan kaidah pada keluaran model, yang dibandingkan antar-model dan antar-tingkat kendali melalui uji statistik. Teori harmoni fungsional dan kaidah Strube digunakan secara deduktif: kaidah diturunkan menjadi variabel, hipotesis, dan definisi operasional sebelum data dikumpulkan.

Berdasarkan jenis rumusan masalah \parencite{sugiyono2019}, pertanyaan penelitian pertama bersifat deskriptif, pertanyaan kedua bersifat komparatif, dan pertanyaan ketiga bersifat asosiatif kausal. Pada pertanyaan kedua, model dipilih oleh peneliti, bukan diberi perlakuan, sehingga perbandingannya bersifat komparatif. Pada pertanyaan ketiga, tingkat kendali diberikan peneliti pada model yang sama, sehingga pengaruhnya diperiksa dengan logika eksperimen, yaitu mencari pengaruh perlakuan tertentu dalam kondisi yang dikendalikan \parencite{sugiyono2019}.

Pemeriksaan contoh partitur pada tahap analisis berfungsi untuk menafsirkan penandaan, bukan sebagai sumber data utama. Pemeriksaan tersebut tidak memiliki pertanyaan, data, dan analisis kualitatif tersendiri, sehingga penelitian ini tidak menggunakan metode campuran (\textit{mixed methods}).

Penelitian berlangsung dalam dua tahap. Tahap pertama membangun dan memvalidasi instrumen evaluasi otomatis berbasis music21. Tahap kedua menerapkan instrumen tersebut untuk menghasilkan temuan komparatif lintas model dan lintas kondisi instruksi.

Perbandingan antar-model diposisikan sebagai studi komparatif, bukan eksperimen terkontrol penuh. Ketiga model tidak sepenuhnya setara dalam kondisi pelatihan: DeepBach dan Coconet merupakan model spesialis yang dilatih pada korpus \textit{chorale} SATB Bach ($\sim$400 karya), sementara NotaGen adalah model generalis yang dilatih pada lebih dari 1,6 juta partitur lintas genre dan dikondisikan ke tugas SATB melalui metadata teks (\texttt{instrumentation: Choir SATB}), bukan melalui \textit{fine-tuning} berbasis genre. Konsekuensinya, perbedaan tingkat kepatuhan yang terukur antar model dapat mencerminkan kombinasi dari (a) perbedaan \textit{bias} induktif arsitektur dan (b) perbedaan skala serta spesialisasi data pelatihan. Desain penelitian ini tidak memungkinkan pemisahan statistik yang bersih antara kedua penjelasan tersebut, karena tidak tersedia versi ketiga model yang dilatih pada data yang identik.

Interpretasi temuan akan membahas kedua penjelasan tersebut dan tidak mengatribusikan seluruh perbedaan skor pada faktor arsitektur semata. Temuan yang dihasilkan berlaku pada model, checkpoint, dan kondisi yang diuji.

Perumusan kaidah dan pemeriksaan partitur tetap melibatkan penilaian peneliti. Posisi peneliti dibentuk oleh latar belakang dan pengalamannya serta dapat memengaruhi proses penelitian \parencite{holmes2020}. Peneliti adalah mahasiswa musik yang mempelajari harmoni Barat, sehingga berposisi sebagai orang dalam (\textit{insider}) terhadap tradisi teori yang dijadikan acuan. Terhadap ketiga model, peneliti berposisi sebagai orang luar (\textit{outsider}) karena tidak terlibat dalam pengembangan dan pelatihannya. Kedua posisi tersebut dapat ditempati bersamaan \parencite{dwyer2009}. Refleksivitas, yaitu pemeriksaan peneliti atas anggapan dan penilaiannya sendiri, juga berguna pada penelitian kuantitatif \parencite{jamieson2023}. Pada penelitian ini, posisi peneliti dapat memengaruhi pemilihan Strube sebagai acuan, perumusan kaidah dan pengecualiannya, anotasi contoh, pemilihan masukan generasi, dan penafsiran hasil. Karena itu, kaidah dan rencana analisis ditetapkan sebelum data utama dikumpulkan, instrumen yang sama diterapkan pada seluruh keluaran, contoh keluaran model dianotasi tanpa informasi model asalnya, dan contoh partitur untuk pembahasan dipilih melalui prosedur yang tercatat. Kaidah Strube mewakili norma penulisan empat suara dalam tradisi harmoni tonal Barat. Penandaan kaidah menunjukkan ketidaksesuaian dengan norma tersebut, bukan penilaian atas nilai musikal keluaran atau atas tradisi musik lain.

\section{Objek dan Subjek Penelitian}

Objek penelitian adalah \textit{output} simbolik (MIDI) yang dihasilkan oleh tiga model AI generatif musik simbolik yang bersifat terbuka (\textit{open source}): DeepBach \parencite{hadjeres2017}, Coconet \parencite{huang2017}, dan NotaGen \parencite{wang2025}. Penelitian ini tidak menjadikan manusia sebagai subjek penelitian. Peneliti dan pengajar harmoni berperan sebagai penganotasi contoh musikal untuk validasi instrumen.

\subsection{Populasi}

Populasi adalah seluruh keluaran yang dapat dihasilkan oleh checkpoint DeepBach, Coconet, dan NotaGen yang ditetapkan dalam protokol, pada setiap kondisi generasi yang dijalankan. Populasi ini tidak dapat dicacah secara lengkap. Karena itu, populasi didefinisikan melalui identitas model, checkpoint, pengaturan generasi, dan masukan yang digunakan.

\subsection{Sampel}

Target jumlah sampel adalah 120 berkas MIDI: 3 model $\times$ 4 kondisi $\times$ 10 sampel per sel eksperimen. Jumlah berkas yang dihasilkan dan jumlah sampel yang memenuhi kriteria analisis merupakan dua hitungan yang berbeda. Setiap percobaan generasi, kegagalan, dan alasan eksklusi dicatat; keluaran yang gagal tidak diganti secara diam-diam.

Masukan generasi (\textit{seed}) dipilih secara purposif berdasarkan kriteria yang ditetapkan dalam protokol. Beberapa keluaran yang dihasilkan dari satu masukan yang sama dapat saling bergantung. Oleh karena itu, jumlah masukan independen dilaporkan terpisah dari jumlah berkas yang dihasilkan, dan unit analisis ditetapkan sebelum pengumpulan data utama.

\section{Metode Pengumpulan Data}

Data dikumpulkan melalui generasi keluaran model dan pengukuran otomatis terhadap keluaran tersebut. Identitas model, checkpoint, versi perangkat lunak, pengaturan generasi, dan masukan setiap run dicatat agar pengukuran dapat diulang.

\subsection{Variabel Penelitian}

Penelitian ini memiliki dua variabel bebas dan satu variabel terikat.

\begin{enumerate}
    \item Variabel bebas pertama ($X_1$): model generatif. Variabel kategorik dengan tiga kategori: DeepBach (LSTM dua arah), Coconet (\textit{deep residual} CNN dengan \textit{Orderless} NADE), dan NotaGen (\textit{hierarchical Transformer} dengan \textit{global self-attention}). Kategori ini menunjukkan model dan checkpoint yang diuji, bukan arsitektur yang terisolasi dari data latih dan prosedur generasi.
    \item Variabel bebas kedua ($X_2$): tingkat kendali (\textit{conditioning level}), yaitu banyaknya arahan musikal yang diberikan kepada model. Variabel kategorik dengan empat tingkat, A sampai D, yang dijelaskan pada Tabel~\ref{tab:kondisi-generasi}.
    \item Variabel terikat ($Y$): tingkat kepatuhan \textit{output} terhadap kaidah Strube. Indikatornya adalah jumlah penandaan kuint sejajar, jumlah penandaan oktaf sejajar, jumlah penandaan kegagalan resolusi nada penuntun, masing-masing beserta penyebutnya, dan \textit{Strube Score} sebagai indeks gabungan. Definisi operasional setiap indikator dijelaskan pada subbab Instrumen Pengukuran.
\end{enumerate}

Pertanyaan penelitian kedua menguji perbedaan $Y$ antar-kategori $X_1$ pada setiap tingkat $X_2$ yang mendukung perbandingan. Pertanyaan penelitian ketiga menguji pengaruh $X_2$ terhadap $Y$ pada setiap kategori $X_1$.

\subsection{Generasi Keluaran Model}

Keempat kondisi dirancang untuk mencerminkan tingkat kendali yang semakin besar. Kondisi dikontrol melalui berkas \texttt{strube\_conditions.json}.

\begin{table}[ht]
\centering
\begin{tabularx}{\textwidth}{|l|l|X|}
\hline
\textbf{Kondisi} & \textbf{Label} & \textbf{Deskripsi} \\ \hline
A & Netral & Generasi tanpa batasan harmoni atau gaya eksplisit \\ \hline
B & Kunci & Generasi dengan konteks kunci/periode musik Barat \\ \hline
C & SATB Parsial & Suara Soprano ditetapkan sebagai batasan; model mengisi Alto, Tenor, Bass \\ \hline
D & SATB Penuh & Soprano dan Bass ditetapkan; model mengisi Alto dan Tenor \\ \hline
\end{tabularx}
\caption{Kondisi generasi}
\label{tab:kondisi-generasi}
\end{table}

DeepBach dan Coconet menerima batasan nada, sementara adaptor NotaGen menggunakan metadata periode, komponis, dan instrumentasi. Pada konfigurasi yang digunakan, kondisi C dan D NotaGen identik, dan kondisi B tidak menentukan C mayor. Dengan demikian, label kondisi tidak menunjukkan intervensi yang setara pada ketiga model. Pelestarian nada yang ditetapkan pada kondisi C dan D diperiksa pada setiap keluaran.

\subsection{Instrumen Pengukuran}

Alat evaluasi dikembangkan menggunakan bahasa Python dan pustaka music21 \parencite{cuthbert2010}. Indikator gabungan variabel terikat adalah indeks yang dinamai \textit{Strube Score} dalam kode. Indeks tersebut mengurangi satu dengan jumlah penandaan kaidah per peristiwa sopran, kemudian membatasi nilainya pada rentang nol hingga satu. Penyebutnya adalah jumlah peristiwa sopran, sehingga beberapa penandaan pada satu waktu dapat berkontribusi secara terpisah terhadap penalti. Jumlah penandaan setiap kaidah juga dicatat secara terpisah.

Kaidah yang diukur bersumber dari Strube \parencite{strube1928}. Rumus pada subbab ini tidak dikutip dari Strube. Peneliti menyusunnya sebagai definisi operasional, yaitu penerjemahan kaidah ke dalam kondisi yang dapat dihitung pada data simbolik. Dalam rumus, $p_{i,t}$ adalah nomor nada MIDI suara $i$ pada peristiwa ke-$t$, interval dihitung dalam semiton, dan kelas nada dinyatakan modulo 12, sesuai representasi nada yang diolah music21 \parencite{cuthbert2010}. $K_{\text{tonic}}$ adalah kelas nada tonika, dan $L_{\text{pc}}$ adalah kelas nada \textit{leading tone}. Rumus ini menjadi spesifikasi yang dicocokkan dengan kode melalui uji perangkat lunak dan dengan penilaian musikal melalui anotasi pada subbab Validitas dan Reliabilitas Instrumen.

Indeks penalti dihitung terhadap $M_{\text{total}}$, yang pada kode saat ini adalah jumlah peristiwa nada/akor sopran dengan nilai minimum satu:

{\small
\begin{multline}
\text{Strube Score} = 1.0 - \min\Biggl(1.0,\; \\
\frac{\displaystyle\sum_{t=1}^{M_{\text{total}}} \left(\sum_{\text{pairs}} \text{Violation}_{\text{parallel}}(t) + \text{Violation}_{\text{leading}}(t)\right)}{M_{\text{total}}}\Biggr)
\end{multline}
}

Deteksi pelanggaran gerak paralel antara dua suara aktif $i$ dan $j$ pada interval target $\delta$ (kuint murni untuk $\delta \equiv 7 \pmod{12}$, oktaf untuk $\delta \equiv 0 \pmod{12}$) diformalkan sebagai:

{\small
\begin{multline}
\text{Violation}_{\text{parallel}}(t) = \mathds{1}\left(|p_{i,t} - p_{j,t}| \equiv \delta \pmod{12}\right) \\
\times \mathds{1}\left(|p_{i,t-1} - p_{j,t-1}| \equiv \delta \pmod{12}\right) \\
\times \mathds{1}\left(\text{sgn}(p_{i,t} - p_{i,t-1}) = \text{sgn}(p_{j,t} - p_{j,t-1})\right) \times \mathds{1}\left((p_{i,t} - p_{i,t-1}) \neq 0\right)
\end{multline}
}

Deteksi kegagalan resolusi nada penuntun pada suara Soprano menerapkan dua filter kontekstual. Pertama, filter modus: pada tangga nada minor, nada ke-7 hanya dianggap \textit{leading tone} jika merupakan \textit{raised seventh} (berjarak satu \textit{semitone} dari tonika, $L_{\text{pc}} \equiv (K_{\text{tonic}} - 1) \pmod{12}$); jika modus terdeteksi sebagai \textit{natural minor} dan nada ke-7 berjarak dua \textit{semitone} dari tonika (\textit{subtonic}), instrumen tidak menandainya sebagai pelanggaran. Kedua, filter arah: pelanggaran hanya dicatat jika \textit{leading tone} bergerak \textbf{naik} ke nada non-tonika ($p_{i,t+1} > p_{i,t}$ dan $p_{i,t+1} \not\equiv K_{\text{tonic}} \pmod{12}$); gerakan turun dari nada ke-7 tidak ditandai sebagai pelanggaran, karena \textit{passing tone} menurun (8-7-6-5) merupakan ornamentasi melodik yang tidak wajib resolusi ke atas. Secara formal:

{\small
\begin{multline}
\text{Violation}_{\text{leading}}(t) = \mathds{1}\left(p_{i,t} \equiv L_{\text{pc}} \pmod{12}\right) \times \mathds{1}\left(p_{i,t+1} > p_{i,t}\right) \\
\times \mathds{1}\left(p_{i,t+1} \not\equiv K_{\text{tonic}} \pmod{12}\right) \\
\times \mathds{1}\left(\text{mode} \neq \text{natural minor} \lor L_{\text{pc}} \equiv (K_{\text{tonic}} - 1) \pmod{12}\right)
\end{multline}
}

Pengecekan gerak paralel dijalankan pada enam pasangan suara SATB: S-A, S-T, S-B, A-T, A-B, T-B. Pengecekan resolusi nada penuntun hanya dijalankan pada sopran. Rumus kuint di atas memakai tujuh semiton modulo 12, sedangkan kode yang ada juga menerima lima semiton; perbedaan tersebut membatasi kesesuaian penandaan kode dengan definisi kuint pada rumus.

Pengecekan \textit{leading tone resolution} dibatasi pada suara Soprano, karena suara tersebut paling konsisten berperan sebagai melodi utama dalam tekstur SATB, sekaligus untuk menghindari kebutuhan deteksi otomatis suara mana yang berperan sebagai pembawa melodi pada \textit{output} AI yang strukturnya bisa jauh dari konvensi SATB klasik. Konsekuensinya, resolusi nada penuntun yang terjadi pada suara Alto, Tenor, atau Bass tidak tercatat oleh instrumen ini. Batasan ini merupakan bagian dari definisi operasional instrumen.

\subsection{Validitas dan Reliabilitas Instrumen}

Sebelum digunakan untuk data utama, instrumen ini menjalani validasi muka (\textit{face validity}) melalui dua pengujian:

\begin{enumerate}
    \item Uji positif: \textit{Chorale} Bach asli dari korpus music21 (BWV 66.6) dievaluasi dan harus menghasilkan skor tinggi, sesuai ekspektasi bahwa \textit{chorale} Bach merupakan contoh penerapan kaidah Strube yang konsisten.
    \item Uji negatif: Harmonisasi yang sengaja dibuat melanggar kaidah Strube (dengan kuint sejajar dan resolusi nada penuntun yang salah secara eksplisit) dievaluasi dan harus menghasilkan skor rendah.
\end{enumerate}

Instrumen yang valid mengukur secara akurat apa yang hendak diukur, sedangkan instrumen yang reliabel memberikan hasil yang konsisten ketika digunakan berulang kali pada objek yang sama \parencite{leedy2018}. Kedua uji di atas merupakan pemeriksaan awal. Kelulusan uji perangkat lunak menunjukkan perilaku pada kasus uji, sedangkan validasi musikal menilai kesesuaian penandaan dengan kaidah dan konteksnya. Validitas instrumen diperiksa dengan membandingkan penandaan algoritmis dengan kumpulan contoh musikal yang dianotasi secara independen. Contoh yang berasal dari keluaran model dianotasi tanpa informasi model dan kondisi asalnya. Setiap anotasi mencatat pasangan suara, konteks tonal, lokasi peristiwa, kaidah, dan alasan penandaan. Kesesuaian penandaan dilaporkan melalui jumlah positif palsu, negatif palsu, \textit{precision}, dan \textit{recall} untuk setiap kaidah.

Reliabilitas instrumen diperiksa melalui pengulangan pengukuran pada berkas yang sama dengan versi instrumen yang sama; pengukuran yang reliabel menghasilkan penandaan yang identik. Apabila anotasi dilakukan oleh lebih dari satu penganotasi, kesepakatan antar-penganotasi dilaporkan sebelum perbedaan anotasi diselesaikan. Perubahan kuantisasi, pemetaan suara, definisi kaidah, atau penyebut skor diperlakukan sebagai versi instrumen baru; hasil versi sebelumnya disimpan.

\subsection{Kontrol Kualitas Sampel}

Tidak semua model generatif menjamin \textit{output} dalam format MIDI empat suara yang terpisah secara eksplisit. Instrumen evaluasi secara otomatis menandai setiap sampel dengan label \texttt{parse\_quality}: \texttt{ok} (empat suara berhasil dipisah), \texttt{padded} (kurang dari empat suara, suara terakhir direplikasi untuk melengkapi), atau \texttt{duplicated\_mono} (hanya satu suara tersedia, direplikasi empat kali). Sampel dengan \texttt{parse\_quality} selain \texttt{ok} dikeluarkan dari analisis utama, dan kegagalan penguraian dilaporkan secara terpisah. Ringkasan yang menggabungkan suara hasil replikasi membatasi interpretasi kepatuhan SATB.

\section{Metode Analisis Data}

Analisis diawali dengan statistik deskriptif: jumlah percobaan, keluaran yang berhasil, kegagalan penguraian, dan sampel yang memenuhi kriteria pada setiap sel; jumlah penandaan setiap kaidah beserta penyebutnya; serta rata-rata dan deviasi standar \textit{Strube Score} per model per kondisi. Statistik deskriptif ini menjawab pertanyaan penelitian pertama.

Selain statistik deskriptif, uji komparatif non-parametrik Kruskal-Wallis direncanakan untuk menguji hipotesis pada Bab II, yaitu apakah perbedaan tingkat kepatuhan antar model (pertanyaan kedua) atau antar-tingkat kendali (pertanyaan ketiga) melampaui ambang signifikansi statistik. Uji non-parametrik dipilih karena ukuran sampel per sel ($n=10$) kecil dan distribusi skor tidak dapat diasumsikan normal tanpa verifikasi terlebih dahulu. Uji Kruskal-Wallis mensyaratkan observasi yang independen dan tidak menunjukkan pasangan kelompok yang berbeda. Uji lanjutan, ukuran efek, dan penanganan perbandingan berganda ditetapkan dalam protokol sebelum data utama dibaca. Apabila keluaran bergantung pada masukan yang sama, analisis menggunakan unit analisis dan uji yang sesuai dengan ketergantungan tersebut.

Hasil disajikan dalam bentuk tabel dan grafik batang kelompok (\textit{grouped bar chart}) yang membandingkan ketiga model di bawah empat kondisi.

Penandaan kaidah merupakan hasil penerapan definisi operasional. Pada korpus Bach maupun keluaran model, penandaan tersebut perlu diperiksa dalam konteks frasa, nada hias, kadens, dan kemungkinan kesalahan penguraian. Detektor tidak menentukan keindahan karya atau maksud komponis. Contoh yang ditandai dipilih melalui prosedur yang tercatat, dianotasi, dan dibahas bersama lokasi birama serta kutipan partiturnya.

\section{Rancangan Alur Penelitian}

Penelitian berlangsung dalam tahap berikut:

\begin{enumerate}
    \item Perumusan kaidah: Penetapan edisi dan halaman kaidah Strube, contoh, pengecualian, dan definisi operasional.
    \item Pengembangan instrumen: Penulisan dan pengujian skrip evaluasi Strube berbasis music21.
    \item Validasi instrumen: Pengujian \textit{face validity} menggunakan \textit{chorale} Bach asli dan harmonisasi yang disengaja melanggar, serta perbandingan penandaan dengan contoh beranotasi.
    \item Uji coba (\textit{pilot}): Generasi dan pemeriksaan sejumlah kecil keluaran dari setiap model dan kondisi untuk menemukan masalah sebelum protokol ditetapkan. Hasil uji coba dipisahkan dari data utama.
    \item Pengumpulan data: Generasi sampel \textit{output} MIDI dari ketiga model di bawah empat kondisi instruksi sesuai protokol yang telah ditetapkan.
    \item Evaluasi otomatis: Seluruh sampel yang memenuhi kriteria dievaluasi menggunakan instrumen yang telah divalidasi; hasil disimpan dalam format CSV terstruktur.
    \item Analisis data: Statistik deskriptif, uji hipotesis, dan pembahasan contoh partitur.
\end{enumerate}

\begin{figure}[htbp]
\centering
\begin{tikzpicture}[node distance=1.45cm, every node/.style={fill=white, font=\footnotesize}, align=center]
  \tikzset{
    startstop/.style={rectangle, rounded corners, minimum width=3.5cm, minimum height=0.7cm, text centered, draw=black},
    process/.style={rectangle, minimum width=3.5cm, minimum height=0.7cm, text centered, draw=black},
    decision/.style={diamond, aspect=2, minimum width=3.5cm, minimum height=0.7cm, text centered, draw=black},
    arrow/.style={thick,->,>=stealth}
  }

  % Nodes
  \node (start) [startstop] {Mulai};
  \node (rule) [process, below of=start] {Perumusan Kaidah Strube \\ \& Definisi Operasional};
  \node (dev) [process, below of=rule] {Pengembangan Skrip \\ Evaluasi Strube (music21)};
  \node (val) [process, below of=dev] {Validasi Instrumen \\ (Uji Positif, Negatif, Anotasi)};
  \node (dec) [decision, below of=val, yshift=-0.1cm] {Apakah Valid?};
  \node (pilot) [process, below of=dec, yshift=-0.3cm] {Uji Coba \\ \& Penetapan Protokol};
  \node (gen) [process, below of=pilot] {Generasi Data Utama};
  \node (eval) [process, below of=gen] {Evaluasi Otomatis \\ \& Simpan CSV};
  \node (anal) [process, below of=eval] {Analisis Data \\ \& Uji Hipotesis};
  \node (stop) [startstop, below of=anal] {Selesai};

  % Connections
  \draw [arrow] (start) -- (rule);
  \draw [arrow] (rule) -- (dev);
  \draw [arrow] (dev) -- (val);
  \draw [arrow] (val) -- (dec);
  \draw [arrow] (dec) -- node[anchor=east] {Ya} (pilot);
  \draw [arrow] (pilot) -- (gen);
  \draw [arrow] (gen) -- (eval);
  \draw [arrow] (eval) -- (anal);
  \draw [arrow] (anal) -- (stop);

  % Feedback loop
  \draw [arrow] (dec.east) -- ++(1.2,0) |- node[anchor=south, pos=0.25] {Tidak} (dev.east);

\end{tikzpicture}
\caption{Rancangan Alur Penelitian}
\label{fig:diagram-alir}
\end{figure}
